\chapter{Introduction}
\label{ch:intro}
Consider four buses travelling between the same two points. Their common road is 2.34 kilometres long. The generated buses arrive after 228, 242, 280 and 308 seconds. A fixed two-minute observation would end before any of them arrived. Figure~\ref{fig:whole-opening} shows the output of the complete-journey sampling program. The endpoint is prescribed, while the model samples the clock and how progress is distributed over it.

\fig{complete-trip.pdf}{1.0}{Four complete generated bus journeys on the same preplanned MoST route. The vertical marker indicates the earlier experimental window. The coloured endpoints in the right panel show the four arrival times. This is the first supported bus example used to check the packaged model.}{whole-opening}

The primary deliverable is a complete timestamped trip for a participant class on a permitted origin-to-destination route. The earlier fixed-window study supplies a controlled comparison of motion before arrival. The road network specifies where a trip can go; it leaves open how long it takes and how the participant moves on the way. We learn a distribution over these possibilities. The group's earlier recognition work motivates learning from spatial and speed representations of the same journey. We study both feedback on sampled trips and direct supervision of the denoiser using these learned representations.

The supplied data consist of two simulation exports without their original route or configuration files. The principal table provides more than four million observations and aligns geometrically with the public Monaco SUMO Traffic network. Reconstruction restricts the population that can support a road-constrained synthesis evaluation. We report that coverage alongside the generated samples.

\section{The synthesis task}
A request specifies a participant class and endpoints on the MoST network, or an externally specified permitted route. A sample is a complete sequence of timestamped positions ending at the destination. We hold the route fixed to compare different motions under identical geometric opportunities. The generator does not receive the observed arrival time. In an earlier controlled task we fixed duration and varied the participant label along a common road. That experiment answers a separate question about the local response of the motion model.

This distinction determines the evaluation. The recorded future is available to the evaluator for population selection and distributional comparisons. The model receives a newly planned path. The empirical comparator receives exactly the same path. A route planner with access to future turns could make a poor motion generator appear geographically accurate. Conversely, an unconstrained coordinate model could reproduce the speed distribution while regularly crossing from one unconnected street to another. The experiments separate these possibilities.

\section{From class recognition to class control}
Previous work at Politecnico di Torino and CNR investigated a classification problem: given a recorded trace, infer the type of traffic participant from an image derived from its ordered positions \cite{salvadori,ferreira}. Our question runs in a different direction. Given a requested participant class, does that representation help select a useful synthetic trace? The original project provides the spatial encoding and image routines. It supplies no trained classifier, so the auxiliary network is trained here on observed trajectories. In the locked experiment it ranks decoded candidates. A later development intervention uses the same score during denoising and tests the effect of changing only the class supplied to the motion network.

There are at least two ways a classifier can succeed in this setting. It can reward plausible class-specific movement, such as a pattern of progression that transfers to another representation. It can also reward spatial or image features preferred by its own architecture. A higher score under the selecting classifier is expected in either case. To distinguish the two explanations, we use an auditor trained on coordinates and train new classifiers with the selected examples. The final measure of training utility is performance on observed held-out participants.

\section{Research questions and contributions}
\begin{enumerate}
\item Can a route-conditioned diffusion model sample a complete, variable-duration journey, and how does its travel-time distribution compare with empirical timing transfer on the same routes?
\item Which parts of class-specific behaviour does the model reproduce? We separate topological feasibility from the temporal arrangement of movement, including sustained cruise and the association between speed and bends.
\item At fixed candidate budget, is selection by a spatial trajectory's class score or by a whole speed sequence's motion-mode score more useful for the delivered class-conditioned dataset?
\item When both spatial and speed teachers learn the same participant labels, can their paired hidden representations improve diffusion synthesis beyond class supervision alone?
\end{enumerate}
The system first makes the learned quantities explicit: the planner supplies a permitted route, the denoiser samples duration and progress, and the decoder produces coordinates through arrival. A matched empirical comparator separates the value of learned timing from the structural benefit of a route. The second contribution is a motion evaluation that can distinguish journeys with the same distance and speed histogram but different temporal ordering. The third is an end-to-end comparison of two feedback policies on the same generated candidates and independently simulated reference traffic. We retain inconclusive and adverse outcomes alongside improvements.

\section{Reading the evidence}
Chapters~\ref{ch:background} and~\ref{ch:data} establish the related work and the provenance of the CNR and SUMO data. Chapters~\ref{ch:generator}--\ref{ch:protocol} explain the model, encoding lineage and evaluation. Chapter~\ref{ch:complete} presents the main complete-journey deliverable and the spatial-versus-speed feedback comparison. Chapters~\ref{ch:results}, \ref{ch:g8} and~\ref{ch:expanded} give the controlled local-motion evidence: empirical transfer, matched class switches, road geometry, data volume and the limits of intermediate scoring. Chapter~\ref{ch:paired} tests paired representation supervision inside diffusion training, with class-centroid and shuffled-pair controls. Chapter~\ref{ch:discussion} brings these results back to the complete-trip task.
